% 探针组：expl3 载入攻坚已修复语义的回归锁（2026-09-12 战役沉淀）
% 每个探针ground truth 一行 P= 期望；abcheck.py 双引擎对拍应全 SAME。
%
% ---- 1. f-展开 protected 停驻（\exp_not:f 语义，fp_set 链前置）----
% tex.web：f 展开展开到首个不可展开 token 为止（protected 宏停驻保留）
\catcode`\@=11 \catcode`\:=11
\let\else:\else \let\fi:\fi
\protected\def\pA{X}
\edef\zB{\detokenize\expandafter{\unexpanded{\pA Y}}}
\message{F-PROTECT=\zB}
%
% ---- 2. edef 内嵌 \exp:w romannumeral 环（fp_parse 展开链前置）----
\edef\zC{\romannumeral-`\0\relax Z}
\message{RN-NEST=[\zC]}
%
% ---- 3. outer 宏作无定界实参首 token：Forbidden + 作业继续 ----
\outer\def\oX{A}
\def\aOne#1{[#1]}
\aOne\oX
\message{OUTER-ARG=CONTINUED}
%
% ---- 4. outer 宏作组实参内 token：同样 Forbidden + 继续 ----
\aOne{\oX}
\message{OUTER-GROUP=CONTINUED}
%
% ---- 5. active char 持 outer 槽：不报 Forbidden（tex.web active 臂）----
\catcode`\^^L=13
\outer\def^^L{\par}
\def\aAct#1{[#1]}
\aAct^^L
\message{ACTIVE-OUTER=NO-FORBIDDEN}
%
% ---- 6. csname 中空格忽略（构名不含空格；prefix 构名依赖）----
\expandafter\def\csname AB C\endcsname{W}
\message{CSNAME-SPACE=[\csname AB C\endcsname]}
%
\end
